\chapter{Analytic Number Theory}\label{analytic-number-theory}
\chapterstatus{planned}

The distribution of prime numbers studied with the tools of analysis.

\noindent\textbf{Prerequisites.} Chapter~\ref{number-theory} (Elementary Number Theory), Chapter~\ref{prime-numbers} (Prime Numbers), Chapter~\ref{complex-analysis} (Complex Analysis) and Chapter~\ref{real-analysis} (Real Analysis).

\noindent\textbf{Planned contents.}
\begin{enumerate}
\item Arithmetic functions and Dirichlet series; Euler products; abscissae of convergence.
\item Averages of arithmetic functions: the hyperbola method, the divisor problem, the Gauss circle problem, squarefree numbers. (Chebyshev's and Mertens' estimates are in Chapter~\ref{prime-numbers}.)
\item The Riemann zeta function: analytic continuation, functional equation, zeros; the prime number theorem (Newman's proof and the classical proof with a zero-free region).
\item The explicit formula and the error term; the Riemann hypothesis and its equivalent forms; Selberg's elementary proof (overview).
\item Dirichlet characters and $L$-functions; non-vanishing of $L(1, \chi)$; Dirichlet's theorem on primes in progressions; the class number formula for quadratic fields (statement).
\item Exercises with solutions.
\end{enumerate}
